\documentclass[12pt]{article}
\usepackage[utf8]{inputenc}
\usepackage{graphicx}
\usepackage{xcolor}
\usepackage{listings}
\usepackage{geometry}
\geometry{a4paper, margin=1in}

\definecolor{codebg}{rgb}{0.95,0.95,0.95}
\lstset{
    backgroundcolor=\color{codebg},
    basicstyle=\ttfamily\small,
    breaklines=true,
    frame=single
}

\title{\textbf{Pharmacy Management Application} \\ Code Explanation for Beginners}
\author{A Simple Guide for Kids and Beginners}
\date{\today}

\begin{document}

\maketitle
\newpage

\tableofcontents
\newpage

\section{Introduction - What is This Project?}

Imagine you want to buy medicine online, just like ordering toys or games! This project is like a digital pharmacy store where people can:
\begin{itemize}
    \item Look at different medicines
    \item Add medicines to their shopping cart
    \item Upload their doctor's prescription
    \item Place orders and buy medicines
\end{itemize}

It's like having a medicine shop on your computer or phone!

\section{Understanding the Main Parts}

Our pharmacy application has two main parts, like two friends working together:

\subsection{The Frontend (What You See) 👀}
This is like the colorful shop window! It's what people see and click on. All the buttons, pictures, and pages are here.

\textbf{Location:} \texttt{frontend/} folder

\subsection{The Backend (The Brain) 🧠}
This is like the shop's storage room and manager! It keeps all the information safe and makes sure everything works correctly.

\textbf{Location:} \texttt{backend/} folder

\section{Important Files Explained Simply}

\subsection{server.js - The Boss of the Shop}

\textbf{Where it lives:} At the main folder (root)

\textbf{What it does:}
Think of \texttt{server.js} as the boss who opens the shop every day! Here's what this boss does:

\begin{itemize}
    \item \textbf{Opens the shop:} When we run the application, this file starts everything
    \item \textbf{Welcomes customers:} It accepts requests from people visiting our website
    \item \textbf{Directs traffic:} It tells the computer where to send different requests
    \item \textbf{Connects to the database:} It makes sure we can talk to our storage of information
    \item \textbf{Handles pictures:} It allows people to upload prescription images
\end{itemize}

\textbf{Simple code explanation:}
\begin{lstlisting}[language=JavaScript]
// This line brings in Express, which helps build websites
const express = require("express");

// This creates our application (like opening a new shop)
const app = express();

// This connects to our database (where we store information)
connectDB();

// These tell the app where to look for different things
app.use("/api/products", productRoutes);
app.use("/api/users", userRoutes);
app.use("/api/orders", orderRoutes);
\end{lstlisting}

\subsection{App.js - The Shop's Front Window}

\textbf{Where it lives:} \texttt{frontend/src/App.js}

\textbf{What it does:}
This is like the main display window of our shop! It decides what page to show when someone clicks on different things.

\begin{itemize}
    \item \textbf{Shows different pages:} Home, Products, Cart, Login, etc.
    \item \textbf{Navigation bar:} The menu at the top where you click to go places
    \item \textbf{Routing:} Like a map that shows which page to go to
\end{itemize}

\textbf{Simple explanation:}
When you click "Home", this file says "Show the Home Page!" When you click "Cart", it says "Show the Cart Page!" It's like a helpful guide in the shop.

\section{Backend Files - The Storage Room}

\subsection{Models Folder - The Information Templates}

\textbf{Where it lives:} \texttt{backend/models/}

Think of models as templates or forms - like when you fill out a form at school!

\subsubsection{User.js - The Customer Information Form}

This file describes what information we save about each customer:
\begin{itemize}
    \item \textbf{Name:} What's your name? (Required - must fill this!)
    \item \textbf{Email:} Where can we email you? (Required and must be unique - no two people can have the same!)
    \item \textbf{Password:} Secret code to login (Required and kept secret!)
    \item \textbf{isAdmin:} Are you a shop manager? (Yes or No)
\end{itemize}

\textbf{Special magic trick:} Before saving the password, it scrambles it so nobody can read it! This is called "encryption" - like writing in secret code.

\begin{lstlisting}[language=JavaScript]
// This is the template for a user
const userSchema = mongoose.Schema({
  name: { type: String, required: true },
  email: { type: String, required: true, unique: true },
  password: { type: String, required: true },
  isAdmin: { type: Boolean, default: false }
});
\end{lstlisting}

\subsubsection{Product.js - The Medicine Information Card}

This file describes each medicine in our shop:
\begin{itemize}
    \item \textbf{unique\_id:} A special number for each medicine (like a barcode!)
    \item \textbf{imgsrc:} A picture of the medicine
    \item \textbf{title:} The name of the medicine
    \item \textbf{price:} How much it costs
    \item \textbf{category:} What type of medicine (for headache, fever, etc.)
    \item \textbf{indication:} What the medicine is used for
    \item \textbf{dosage:} How much to take
    \item \textbf{sideEffects:} Things to watch out for
    \item \textbf{countInStock:} How many we have in the shop
\end{itemize}

\subsection{Routes Folder - The Direction Signs}

\textbf{Where it lives:} \texttt{backend/routes/}

Routes are like signs in a big building that say "Bathroom this way →" or "Exit that way ←"

\subsubsection{userRoutes.js - The Customer Service Desk}

This file handles all requests about users:
\begin{itemize}
    \item \texttt{POST /api/users} - Register a new customer
    \item \texttt{POST /api/users/login} - Let a customer login
    \item \texttt{GET /api/users/profile} - Show customer their profile
    \item \texttt{PUT /api/users/profile} - Let customer update their profile
\end{itemize}

Think of it like different windows at a bank - one for new accounts, one for checking your balance, etc.

\subsubsection{productRoutes.js - The Medicine Information Desk}

This handles requests about products:
\begin{itemize}
    \item Get all medicines
    \item Get one specific medicine
    \item Add a new medicine (only for admin/shop owner)
    \item Update medicine information
    \item Delete a medicine
\end{itemize}

\subsection{Controllers Folder - The Workers}

\textbf{Where it lives:} \texttt{backend/controller/}

Controllers are like the actual workers who do the job! When someone asks for something, the controller makes it happen.

\subsubsection{userControllers.js - The Customer Service Worker}

This file contains functions that actually do the work:

\textbf{registerUser:} Creates a new customer account - like filling out forms and creating a membership card

\textbf{authUser:} Checks if the email and password are correct - like checking if you have the right key to open a door

\textbf{getUserProfile:} Gets the customer's information to show them - like looking up your homework from a file

\subsection{Middleware Folder - The Security Guard}

\textbf{Where it lives:} \texttt{backend/middleware/authMiddleware.js}

Middleware is like a security guard at the entrance! Before letting you do certain things, it checks:
\begin{itemize}
    \item Are you logged in?
    \item Do you have permission to do this?
    \item Is your token (like an ID card) valid?
\end{itemize}

If you pass the checks, you can continue. If not, the guard says "Sorry, you can't enter!"

\section{Frontend Files - The Shop Display}

\subsection{Screens Folder - Different Shop Pages}

\textbf{Where it lives:} \texttt{frontend/src/screens/}

Each file here is like a different room in our shop:

\subsubsection{HomeScreen.js - The Welcome Room}
This is the first page people see - like the entrance of a store! It shows:
\begin{itemize}
    \item Welcome messages
    \item Featured products
    \item Navigation to other parts
\end{itemize}

\subsubsection{AllProductsScreen.js - The Medicine Shelf}
This shows all the medicines available, like walking down the aisles in a pharmacy.

\subsubsection{CartScreen.js - The Shopping Basket}
This shows what medicines you picked - like your shopping cart! You can:
\begin{itemize}
    \item See what you selected
    \item Remove items
    \item Change quantities
    \item Proceed to checkout
\end{itemize}

\subsubsection{LoginScreen.js - The Member Login Door}
This is where customers enter their email and password to login.

\subsubsection{ProfileScreen.js - Your Personal Locker}
This shows your personal information and order history.

\subsubsection{PrescriptionScreen.js - The Doctor's Note Upload}
This is where customers can upload their doctor's prescription picture.

\subsubsection{AdminScreen.js - The Manager's Office}
Only the shop manager can enter here! They can:
\begin{itemize}
    \item Add new medicines
    \item Remove medicines
    \item Update prices
    \item See all orders
\end{itemize}

\subsection{Redux Folder - The Shop's Memory System}

\textbf{Where it lives:} \texttt{frontend/src/redux/}

Redux is like the shop's memory notebook! It remembers:
\begin{itemize}
    \item What's in your cart
    \item Who is logged in
    \item What products are available
\end{itemize}

\subsubsection{Actions - The Commands}
Actions are like giving commands: "Add this to cart!", "Remove that!", "Login now!"

\subsubsection{Reducers - The Memory Keepers}
Reducers listen to the commands and update the memory. If the action says "Add to cart", the reducer writes it down in the memory.

\subsubsection{store.js - The Master Memory Book}
This is the main memory book that keeps track of everything!

\subsection{Components Folder - The Reusable Parts}

\textbf{Where it lives:} \texttt{frontend/src/components/}

Components are like LEGO blocks! Instead of building the same thing again and again, we create one block and use it everywhere.

\subsubsection{Navbar.js - The Top Menu Bar}
The navigation bar at the top of every page - like the menu at the top of a restaurant.

\subsubsection{Product.js - The Medicine Card}
This is a small card that shows one medicine. We use this same card for every medicine! Just like using the same picture frame for different photos.

\subsubsection{CartItem.js - The Basket Item Card}
A small card that shows one item in your cart.

\subsubsection{Footer.js - The Bottom Information}
The footer at the bottom of every page with contact info and links.

\section{Configuration Files}

\subsection{db.js - The Database Connection}

\textbf{Where it lives:} \texttt{backend/config/db.js}

This tiny but important file connects our application to MongoDB (our big storage box of information).

\textbf{What it does:}
\begin{lstlisting}[language=JavaScript]
// Try to connect to the database
await mongoose.connect(process.env.MONGO_URI);
console.log("MongoDB connection SUCCESS");
\end{lstlisting}

It's like plugging in a USB drive to your computer - making the connection so data can flow!

\subsection{package.json Files - The Shopping Lists}

There are two \texttt{package.json} files - one for frontend, one for backend.

Think of these as shopping lists! They tell the computer:
\begin{itemize}
    \item What tools we need (called dependencies)
    \item What commands we can run
    \item Information about our project
\end{itemize}

For example:
\begin{itemize}
    \item We need "express" - to build the website
    \item We need "react" - to make the pages interactive
    \item We need "mongoose" - to talk to the database
\end{itemize}

\section{How It All Works Together}

Let's say someone wants to buy medicine. Here's the journey:

\begin{enumerate}
    \item \textbf{Customer visits website:} They type the website address and App.js shows them the HomePage
    
    \item \textbf{Browsing products:} They click "Products" and AllProductsScreen.js shows all medicines
    
    \item \textbf{Getting data:} The frontend asks the backend "Give me all products!" by making a request to \texttt{/api/products}
    
    \item \textbf{Backend responds:} The server.js receives the request, routes it to productRoutes.js, which calls productControllers.js, which asks the database for all products
    
    \item \textbf{Sending back data:} The database gives the products to the controller, which sends them back to the frontend
    
    \item \textbf{Showing the products:} The frontend receives the data and displays it on the page using Product component cards
    
    \item \textbf{Adding to cart:} Customer clicks "Add to Cart" and Redux stores it in memory
    
    \item \textbf{Checking out:} Customer goes to cart, fills in details, and places order
    
    \item \textbf{Saving order:} Backend receives order and saves it to database
\end{enumerate}

It's like a relay race where each part passes the baton to the next!

\section{Summary - The Big Picture}

\textbf{Think of it like a real pharmacy:}
\begin{itemize}
    \item \textbf{Frontend = Shop Display:} The colorful store people walk into
    \item \textbf{Backend = Storage Room \& Manager:} Where products are kept and managed
    \item \textbf{Database = Big Filing Cabinet:} Where all information is stored safely
    \item \textbf{Routes = Direction Signs:} Telling where to go for different things
    \item \textbf{Controllers = Workers:} Actually doing the tasks
    \item \textbf{Models = Forms \& Templates:} Standard formats for information
    \item \textbf{Middleware = Security Guards:} Making sure only authorized people can do certain things
    \item \textbf{Components = LEGO Blocks:} Reusable pieces we use everywhere
    \item \textbf{Redux = Memory System:} Remembering what's happening
\end{itemize}

\textbf{Remember:} Just like a real shop needs many people and systems working together, our pharmacy application has many files and folders working as a team!

\end{document}
